% p1-appC-phase-transition

\section{P1 Appendix C: Phase Transition Structure (Conjecture)}

Appendix C: Phase Transition Structure (Conjecture)

       Figure (fig-p1-appC-phase-transition). Phase-transition triptych --- Order parameter / Critical
                                       point / MDL evidence bars.

  C.1 Coverage Functional and Order Parameter

        Conjecture 1 (Pellis Symbolic Phase Transition). There exists a critical
        complexity scale C\textasciicircum{} > 0 such that the coverage functional satisfies:*

  rho(C, $\varepsilon$) = O(e-alpha(C\textasciicircum{} - C)), \& C < C\textasciicircum{}, 1 - O(e-beta(C - C\textasciicircum{})), \& C > C\textasciicircum{}, C.1

  for constants alpha, beta > 0 determined by the effective metric entropy of T in
  log-coordinates.

  This conjecture is motivated by the analogy with percolation phase transitions in
  random lattice models, where a sharp onset of coverage occurs when the lattice
  spacing crosses the scale of inter-point distances. It is labeled as a conjecture and has
  no formal proof in this paper.

  C.2 Characterization of the Critical Point
  The critical scale C\textasciicircum{} is characterized by the condition that the effective lattice spacing
  of S(C\textasciicircum{}) in log-space equals the typical inter-point distance of log T:

  deltaeff(C\textasciicircum{}*) $\sim$ E[d(log Xi, log Xj)]. C.2

  The susceptibility chi(C) := $\partial$ rho / $\partial$ C is conjectured to achieve a global maximum near
  C\textasciicircum{}*, with width Delta C $\sim$ |T|-1/2.

  C.3 Relation to MDL Evidence
  Under the conjecture,

  DeltaMDL(C\textasciicircum{}) approx 0, (d DeltaMDL)/(dC)|C = C\textasciicircum{} gg 1. C.3

  The critical point C\textasciicircum{}* corresponds to the grammar complexity at which the structured
  grammar first achieves a net coding advantage over the null model.

